\section{Conclusions and Future Outlook}
\label{sec:conclusions}

In this work, we have established a mathematically rigorous, singularity-free framework for stable boundary layer (SBL) modeling by unifying Geometric Singular Perturbation Theory (GSPT) with the differential geometry of vertical profiles \cite{nieuwstadt1984}. By re-parameterizing Nieuwstadt local scaling via the inverse Obukhov length profile $\chi(z) \equiv 1/L(z)$, we eliminated numerical singularities near neutral conditions ($H \to 0, L \to -\infty$) and proved that visual "knees" observed in physical vertical profiles are frequently **"Fold Illusions"** \cite{sheba2002}. These spatial structures are driven by coordinate-induced mapping curvature ($Ri_{g,zz} = \mathcal{C}_{\text{constitutive}} + \mathcal{C}_{\text{mapping}}$) under active vertical flux divergence, rather than genuine structural failures of turbulent transport closures.

Dynamically, we mapped SBL turbulence collapse as a slow-passage bifurcation across a folded, three-branch critical manifold $\mathcal{C}_0$. Transition to the decoupled, very stable boundary layer (VSBL) regime occurs at a saddle-node fold locus where the fast eigenvalue vanishes ($\lambda_f \to 0$), destroying Fenichel normal hyperbolicity \cite{fenichel1979}. To stabilize discrete NWP and SCM solvers across these delicate regime boundaries, we introduced three structural formulations:
\begin{enumerate}
    \item \textbf{Adaptive Equidistribution Grid Stretching ($\Delta z \propto 1/|\zeta_{zz}|$):} Concentrates vertical nodes dynamically around migrating low-level jet noses and high-curvature zones to resolve steep coordinate gradients, achieving quadratic spatial convergence $\mathcal{O}(N_z^{-2})$ while remaining strictly bounded between $\Delta z_{\text{min}}$ and $\Delta z_{\text{max}}$ to preserve CFL constraints \cite{gabls3_2013}.
    \item \textbf{Fenichel Invariant Manifold Slaving ($u_2 = \Phi(u_1, S)$):} Replaces prognostic fast-mode ODEs with regularized algebraic slow-manifold parameterizations, eliminating spurious cubic Galerkin blowups and establishing globally bounded solver stability \cite{fenichel1979}.
    \item \textbf{Asymptotic Boundary Regularity Constraints ($b - \frac{3}{2}a \ge 2$):} Eliminates coordinate singularities at the SBL top by forcing turbulent momentum ($\tau$) and heat ($H$) fluxes to decay quadratically ($a=2$) and quintically ($b=5$) near inversion boundaries.
\end{enumerate}

These theoretical and numerical formulations were validated against CASES-99 tower observations and GABLS3 single-column benchmarks \cite{cases99, gabls3_2013}, confirming that our curvature error-partitioning schema and gradient-masking weighting function $w(z)$ prevent unphysical parameter drift and maintain numerical stability under non-vanishing CFL time-step limits.

\subsection{Future Research Directions}

Three key avenues emerge to expand upon this GSPT framework:
\begin{enumerate}
    \item \textbf{Doppler Lidar Profile Integration:} Extending our pre-diagnostic Tikhonov splines to Doppler Lidar remote-sensing profiles to account for range-dependent signal-to-noise ratios (SNR), where measurement uncertainty $\delta(z)$ increases with altitude, preserving inversion-layer gradients.
    \item \textbf{WSINDy Dynamical Discovery:} Feeding regularized, non-singular coordinate fields directly into Weak-form Sparse Identification of Non-linear Dynamics (WSINDy) algorithms to discover candidate governing PDEs for SBL transitions from field data, directly testing our fast-slow prognostic hypotheses.
    \item \textbf{Machine Learning Regime Classification:} Training explainable classifiers on GSPT coordinate trajectories—specifically tracking the migration of the fold locus $\zeta_z(z_*) = 0$ in height-time space—to detect nocturnal transitions and decoupling events in real time.
\end{enumerate}

Embedding these continuous geometric constraints into operational NWP grid generation and numerical solver architectures provides a foundation for physically consistent, grid-independent stable boundary layer parameterizations.